% !TEX root = ../monograph.tex
% Current defect and uniform transfer bounds.
\section{Current conservation and uniform transfer bounds}
\label{ch17:quadratic-current}

The potentially dangerous quadratic defect has an exact cancellation
that should be imposed before estimating its norm. We prove preservation
of that cancellation for the linearization of the critical configuration. This removes
an additive interpretation of the order-two matching error. We then control the nonlinear
jets of the critical configuration and the current itself.

Fix a stage and an admissible action. Let $\mathcal E_b$ be a
gauge-invariant current defect with a base point, with zero constant and linear
field coefficients. Its Hessian at zero has an exponentially local
joint Fourier symbol $E_b(p,q)$, with two independent field momenta.
Gauge invariance gives a Ward identity on each leg. We say that its
quadratic defect has order at least four when
\begin{equation}
                    E_b(p,q)=O((|p|+|q|)^4)
                              \quad\hbox{near }(0,0).
 \label{ch17:defect-four}
\end{equation}
Equivalently, every Taylor coefficient of $E_b$ of total degree less
than four in $(p,q)$ vanishes. This order of vanishing in momentum
implies the cancellations for pairs of affine fields and for an affine
field paired with a quadratic field that are used in the current estimate.

For the initial Wilson current, \eqref{ch17:defect-four} follows from
the leading continuum identity. Its linearized curvature is $dA$;
the leading quadratic action is the curvature square; and the centered
difference $\nabla$ has linear symbol $ip$. The quadratic Taylor
coefficient of
$\sum_a\nabla_aJ_{ab}-DS[Y_b]$ therefore has leading spatial
momentum degree three. At that degree it is exactly
\[
 \sum_a\partial_a\left(\tfrac14\delta_{ab}\sum_{c,d}\|F_{cd}\|_{\gLie}^2
                         -\sum_c \langle F_{ac},F_{bc}\rangle_{\gLie}\right)
       +\sum_{a,c} \langle F_{bc},\partial_aF_{ac}\rangle_{\gLie}=0.
\]
The identity uses the linear Bianchi equation. Lower spatial degrees vanish because each curvature
contains a difference and the action variation contains one more.
All remaining lattice Taylor terms have degree at least four.
The same reasoning applies to any initial path action with the
normalized Taylor polynomial of scaling dimension four already established in
Chapter~\ref{ch07:matching}.

\subsection{Preservation of the order of vanishing}

The first derivative of the nonlinear critical-configuration map, from coarse
to fine link variables with block ratio $L$, is the forest representative
$W_F$ of \eqref{ch05:stationary-tangent}; its higher derivatives
cannot enter the quadratic pullback because the defect has zero first
field derivative. Since $W_F-W_a$ is a fine gradient, the Ward identities
of $E_b$ allow us to compute with the smooth representative $W=W_a$.

For the exact Taylor-order statement we pass once more to a gauge-equivalent representative, the linearization $W_s$ built with the linear averaging operator $M$ (Section~\ref{ch05:actual-tangent}), and test Hessians on polynomial fields.\mgubi{Issue 8: polynomial fields live on $\ZZ^4$, the critical configuration of Chapter~\ref{ch05:gaussian} on a torus. One sentence is missing: the Hessian on the torus is the periodization of the one on $\ZZ^4$. Vanishing on a constant leg (used below) is false on a torus. Counterparts: existence/\allowbreak quadratic-pole-transport.\allowbreak tex; v2 I-section.\allowbreak tex.} By \eqref{ch05:corrected-transport}, $W_a-W_s=-d_{0,f}ERW_s$; since $E$ is bounded and $RW_s$ has a bounded coarse kernel, each column of this difference is the gradient of a bounded fine potential. Let $F$ be a gauge-invariant local function of the fine field with base point $o$ and zero first field derivative, and write $B(a,a')=D^2F(0)[a,a']$, whose kernel decays exponentially in the distances of both legs from $o$.
The Ward identity on each leg, $B(d_{0,f}\phi,a')=B(a',d_{0,f}\phi)=0$, holds for finitely supported $\phi$ and, by this decay, also for bounded $\phi$ and fields $a'$ of polynomial growth: replace $\phi$ by its restriction to a large ball about $o$. Call a one-form with values in $\gLie$ a polynomial field of degree $r$ if its components in a basis of $\gLie$ are polynomials of degree at most $r$ in the coordinates. $B$ converges absolutely on pairs of polynomial fields, and the Taylor coefficient of multidegree $(\alpha,\alpha')$ at zero of the joint symbol of $B$, expanded about $o$, is a nonzero multiple of the value of $B$ on the fields $(x-o)^\alpha e$ and $(x'-o)^{\alpha'}e'$ with $e,e'$ constant. The symbol being analytic, \eqref{ch17:defect-four} for $F$ is therefore equivalent to
\begin{equation}
 B(P,P')=0\qquad\hbox{for polynomial fields $P,P'$ of degrees $r,r'$ with }r+r'\le3.
 \label{ch17:polynomial-four}
\end{equation}
This condition does not depend on the center of the coordinates, since a translate of a polynomial has the same degree. Hence it holds for every fine translate of $E_b$, and so for $F=K_0E_b$, where $K_0$ is the finite combination of fine translates of the base point from Lemma~\ref{ch17:block-identities}; this $F$ is again gauge invariant with zero first field derivative, and its kernel decays exponentially because $K_0$ has range $CL$.
Let $G=\mathcal IF$ be its composition with the critical configuration. For finitely supported coarse fields $a,a'$ the preceding paragraph gives $D^2G(0)[a,a']=B(W_Fa,W_Fa')=B(W_aa,W_aa')$, where the potential of the fine gradient $W_Fa-W_aa$ is bounded because it is a sum of values of $W_aa$ along the spanning tree of one block; and the Ward identity applied to the bounded potentials above gives $D^2G(0)[a,a']=B(W_sa,W_sa')$. The kernel of $W_s$ is $O(L^{-1})$ with exponential coarse decay (Section~\ref{ch05:transports}).
Hence the coarse Hessian kernel of $G$ decays exponentially in the distances of both legs from the coarse block of $o$, its joint symbol is analytic, and for coarse polynomial fields $P,P'$ absolute convergence gives $D^2G(0)[P,P']=B(W_sP,W_sP')$. As for $F$, \eqref{ch17:defect-four} for $G$ is equivalent to the vanishing of these values when $r+r'\le3$.
By \eqref{ch17:polynomial-response}, whose derivation in the proof of Proposition~\ref{ch17:spatial-moments} uses only the formulas of Chapter~\ref{ch05:gaussian}, $W_s$ maps a coarse polynomial field $P$ of degree $r<j$ to a fine polynomial field of degree at most $r$; its homogeneous part of degree $r$ is $L^{-r-1}$ times that of $P$, read in fine coordinates. Here $j$ is the smoothing order of the averaging weight, and $j\ge5$, since \eqref{ch05:actual-bridge} is used below with three fine differences.\mgubi{Issue 8: $j\ge5$ is stronger than the standing $j\ge4$ of Chapter~\ref{ch05:gaussian} and is not entered in the schedule of Chapter~\ref{ch10:trajectory}; nor are the requirements on the smoothing order in Section~\ref{ch11:marked-step} and in the proof of Proposition~\ref{ch17:spatial-moments}.}
If $r+r'\le3$, the two images therefore have degrees with sum at most three, and \eqref{ch17:polynomial-four} for $F$ gives $D^2G(0)[P,P']=0$. Thus $G$, and the pulled-back defect $L^5G$, satisfy \eqref{ch17:defect-four} on the coarse lattice.
The same computation gives the exact rescaling of the first coefficients that need not vanish. Suppose that $B$ vanishes on all pairs of polynomial fields with $r+r'<n$, and let $P,P'$ be homogeneous of degrees $r,r'$ with $r+r'=n<j$. Then the parts of $W_sP$ and $W_sP'$ of lower degree contribute nothing, and $D^2G(0)[P,P']=L^{-n-2}B(P,P')$, with $P,P'$ read in fine coordinates on the right.
No alias decomposition of the linearization and no projection off the fine gradient direction is used: fine gradients are removed by the Ward identity before any polynomial is inserted, and only the polynomial moments of $W_s$ enter.
No bound uniform in $L$, and no inverse of a kinetic symbol at a nonzero alias, is
needed for this exact Taylor-order statement.

The base-point average $K_0$, the average over translates of the base point, from
Lemma~\ref{ch17:block-identities}, is performed before composition with the critical
configuration. Its symbol has a double zero at every nonzero
alias: at least one coordinate of $l$ is nonzero, and the factor
$M_a^2G_a(s_a^L)$ has two factors of the average $M_a$ there.
The argument above does not use this double zero. Of $K_0$ it uses only
that it is a finite combination of fine translates of the base point; the
behaviour of the linearization enters only through \eqref{ch17:polynomial-response}.

Thus the order of vanishing is preserved by an exact computation. The
pulled-back Hessian, evaluated on pairs of coarse polynomial fields of total degree
at most three, equals the Hessian of $K_0\mathcal E_b$ evaluated on their images
under $W_s$. These images are fine polynomial fields of no larger total degree,
on which that Hessian vanishes by \eqref{ch17:polynomial-four}. No estimate,
no limit and no bound uniform in $L$ enters. Hence $L^5\mathcal I(K_0\mathcal E_b)$
satisfies \eqref{ch17:defect-four} on the coarse lattice. This is an
exact identity and holds on the same history.

In particular the quadratic jet is not merely close to having the
required affine cancellations: it has them exactly, provided the
transfer of the current uses the exact pulled-back defect
$\mathcal E'_b=L^5\mathcal I(K_0\mathcal E_b)$ through the relevant
field order. If an implementation redefines that defect by action
matching, equality of these quadratic jets must be retained; a bounded
additive matching error would not preserve this argument.

\subsection{A homogeneous bound for the quadratic channel}

The analytic division argument of Section~\ref{ch07:division} can be
applied to each of the two independent momentum legs of $E_b$.
The Ward identities produce one curvature difference on each leg.
Their homotopy divisions lower Taylor order by one each, so
\eqref{ch17:defect-four} leaves a coefficient of total order at least
two. Coordinatewise Taylor division in the remaining eight variables
factors two additional differences, distributed between the legs.
On a fixed smaller coefficient width these operations are bounded by
constants depending only on the stage and the decrease of the width. Thus $E_b$
is a sum of terms with either one and three differences, or two and
two differences, on its two field legs. The coefficient sum is bounded
by $C_j\|E_b\|$ in the original weighted norm.

Recall that $W=W_a$ is the smooth representative; the quadratic Ward
identities annihilate its gradient difference from $W_F$.
Its fine-difference bounds are
\[
                  \|\nabla_f^rW_a\|\le C_{j,r}L^{-1-r},
                                  \qquad r\le3.
\]
Consequently either allocation of four differences contributes at most
$C_jL^{-6}$ for the two legs. The base-point average $K_0$ has bounded
absolute coefficient sum, not a factor $L^4$; it is an average rather
than an unnormalized sum. Translated coefficient weights are absorbed
by a fixed narrower coarse width, exactly as in the inherited quadratic
action estimate of Chapter~\ref{ch07:matching}. Multiplication by the
current normalization $L^5$ proves
\begin{equation}
             \|D^2\mathcal E'_b(0)\|_{\rm out}
                  \le \frac{C_j}{L}\,
                            \|D^2\mathcal E_b(0)\|_{\rm in}.
 \label{ch17:quadratic-current-contraction}
\end{equation}
This is homogeneous in the defect. No constant inhomogeneous term has
been introduced. Dividing the coefficients by $(\log\widehat\beta)^2$
only adds the adjacent logarithmic ratio, which is bounded and tends
to one on a fixed-stage tail. The selected block ratio makes this
quadratic channel a strict contraction.

\subsection{Ratio-dependent bounds for the nonlinear terms}

The nonlinear jets of the critical configuration enter only below the diagonal in field
order. Their constants may depend on $L$, because the lower bound on the reference coupling is
chosen after the ratio. What must remain uniform before $L$ is the
coarse coefficient width on which these constants are measured.
We obtain that width from the gauge-fixed inverse of
Chapter~\ref{ch07:matching}.

Write $z(v)$ for the fluctuation coordinates of the critical configuration, and
$\mathscr S(v,z)$ for the same joint classical action as there.
Let $T u=(u,Dz(0)u)$, and let $\mathscr H$ be its Hessian in the fluctuation coordinates.
On the slice of the block axial gauge implicit differentiation gives
\begin{align}
 \mathscr H z_2[u,v]
    &=-D^3\mathscr S[T u,T v,\,\cdot\,],\notag\\
 \mathscr H z_3[u,v,w]
    &=-D^4\mathscr S[T u,T v,T w,\,\cdot\,]
      -\sum_{\{ij\}\{k\}}
             D^3\mathscr S[(0,z_2[u_i,u_j]),T u_k,\,\cdot\,].
 \label{ch17:stationary-higher-jets}
\end{align}
There are three terms in the sum. The right sides are fluctuation
covectors, with all mixed derivatives in the averaging operation and in the local gauge transformation included.

These formulas have period-independent bounds $C_{j,L}$ on one
fixed smaller coarse Fourier strip. To check the width statement in this zero-background computation, take the forest of
Section~\ref{ch05:actual-coordinates} to consist of the coarse translates of one reference tree. This fixes a free choice for the present computation only; the average over forests of Section~\ref{ch09:forest-average} is not used here, and independence of the forest is not claimed for the forest-dependent jets $z_2,z_3$ themselves. The bounds are used below only for pullbacks of gauge-invariant functions along the critical-configuration map; that these do not depend on the forest (the slices of the block axial gauge for two forests differ by a gauge transformation equal to one at the block roots) is used here and is not proved in this chapter.\mgubi{Issue 7: v3 has no such lemma. Lean has it near the zero configuration (YM.\allowbreak Work.\allowbreak F05.\allowbreak answer, from graph\_switch, fine\_switch, Forest.\allowbreak fold\_root; the proofs of the first two were not read). For the independent forests of Chapter~\ref{ch09:renormalization}, $\Pi_F$ does not commute with coarse translations.} Let $\Pi_F$ be projection onto the slice of the block axial gauge along the joint gauge
directions; it then commutes with coarse translations. At zero it is computed by finite sums along the block tree
and the local gauge transformation at the root; its coefficients and those of $\Pi_F^*$ have norm
bounded by a finite $C_{j,L}$, but bounded range in coarse blocks.
The dual $\Pi_F^*$ extends a slice covector by zero on gauge directions.
For such a covector the gauge-fixed inverse $\mathscr C^{-1}$ of
Chapter~\ref{ch07:matching} solves the Hessian equation modulo gauge.
Thus the slice inverse is
\[
                       \Pi_F\mathscr C^{-1}\Pi_F^*.
\]
The gauge-fixed inverse has a common positive coarse width chosen before
$L$; multiplication by the finite-range projections does not destroy it.
The path vertices and derivatives of the averaging operation also have a common coarse
width, with finite $C_{j,L}$ norms. Restrict the external strips by a
fixed factor so all intermediate momentum sums lie in this inverse
strip. Formula~\eqref{ch17:stationary-higher-jets} then proves the
asserted bounds for $z_2,z_3$. It does not assert bounds uniform in $L$
for their norms.

For every gauge-invariant function with a base point, evaluate this critical
configuration in a radial gauge at the base point. At first order, the gauge-corrected
link at fine displacement $x$ is a sum of the curvature of the linearization over
coordinate rectangles of total area at most $C(1+|x|)$. The smooth
curvature bound for $W_a$ gives
\[
 |W^{\rm rad}(x,y)|
 \le C_jL^{-2}(1+|x|)
                   e^{C|x|/L}e^{-c_jd(0,y)}.
\]
The polynomial factor and fine exponential are absorbed by the weighted
local coefficient norm. A common strip reduction handles the coarse
weight. Higher derivatives of this gauge change are finite products
along the same contours; the fixed-order factors are polynomial in contour
length. Their period-independent norms are therefore still finite
$C_{j,L}$ on the same smaller coarse width.

The raw higher coefficients $D^3\mathcal E,D^4\mathcal E$ need not
separately satisfy linear Ward identities. We use gauge invariance of
the complete function to make this radial substitution first; its
nonlinear derivatives multiply only lower field derivatives of the
defect.

Now differentiate the exact pullback $\mathcal E'=L^5\mathcal I(K_0\mathcal E)$.
In field order $r$ the term involving $D^r\mathcal E$ has only $r$
factors of the linearization of the critical configuration. For $r=3,4$, radial gauge bounds it by
$C_jL^{5-2r}\|D^r\mathcal E\|$. The quadratic diagonal has the
stronger estimate \eqref{ch17:quadratic-current-contraction}.
Every other chain-rule term involves $D^s\mathcal E$ with $s<r$
and higher jets of the critical configuration or of the gauge change. Hence, for $r=2,3,4$,
\begin{equation}
 e'_r\le a_{j,r}(L)e_r+\sum_{2\le s<r}C_{j,L,r,s}e_s,
 \qquad
 a_{j,2}(L),a_{j,3}(L)\le C_j/L,
 \quad a_{j,4}(L)\le C_j/L^3.
 \label{ch17:unnormalized-defect}
\end{equation}
Here $e_r$ is the weighted norm of the field derivative of order $r$.
The estimates are homogeneous in the defect, because every term in the
composition contains one of its derivatives. There is no action-independent
inhomogeneous term on the right.

Put $\bar e_{m,r}=e_{m,r}/(\log\widehat\beta_m)^r$. For a fine-to-coarse
step from $m+1$ to $m$, the coefficient below the diagonal becomes
\[
 C_{j,L,r,s}
       \frac{(\log\widehat\beta_{m+1})^s}
            {(\log\widehat\beta_m)^r},\qquad s<r.
\]
At fixed $j,L$ this tends to zero as the lower bound on the reference coupling increases.
The diagonal logarithmic ratio tends to one. Choose $L$ first so the
three diagonal coefficients are at most $1/8$, then choose the lower bound on the reference coupling
so each row sum, including the below-diagonal terms, is at most $1/2$.
The finitely many adjacent-stage cases are included before each choice.
This proves a homogeneous contraction of the normalized defect vector.
Starting with the Wilson defect, whose normalized norm is bounded, its
value after $n$ transfers is at most $C2^{-n}$. At a fixed scale
the logarithmic normalization is fixed, so the unnormalized
defect derivatives tend to zero as well.

\subsection{Why action matching introduces no inhomogeneous term}

We specify the map on vector fields, including its sign. At the critical
configuration and coarse local gauge transformation, let $\widetilde Q(U,h)$ be the
averaging operation with its chosen local gauge transformation. Its left-trivialized differential is
\[
 (R_f\eta)_e=\left.\frac d{ds}\right|_0
 \log\bigl(\widetilde Q(e^{s\eta}U,h)_e
                              \widetilde Q(U,h)_e^{-1}\bigr).
\]
Let $b_{y,f}$ be the fixed contour coefficients of the local gauge transformation, defined by
$\sum_f b_{y,f}\eta_f=\tau_y(0;\eta)-\sum_rw_r^{\rm av}\tau_y(r;\eta)$,
where $\tau$ is the signed path average in
\eqref{ch05:frame-linearization}. With the fixed contour $\gamma_{y,f}$
from coarse root $Ly$ to the initial vertex of the fine link $f$, set
\[
 (R_0\eta)(y)=\sum_f b_{y,f}
             \operatorname{Ad}(h(y)U(\gamma_{y,f}))\eta_f,
 \qquad (d_V\phi)_e=\phi(e^-)-\operatorname{Ad}(V_e)\phi(e^+).
\]
At coarse base point zero the transferred coarse vector field is
\begin{equation}
 Y'_b=L^5\sum_xK_0(x)(R_f+d_VR_0)
                              Y_b^x(U_{\rm stat}(V)).
 \label{ch17:actual-current-map}
\end{equation}
Translations give the other base points. The sum is over fine base points;
the output is a vector on coarse links. At the identity $R_0=R$ of
\eqref{ch05:actual-Q}, $R_f=Q=M+d_{0,c}R$, and
$d_V=-d_{0,c}$. In this convention the plus sign in
\eqref{ch17:actual-current-map} cancels the root correction exactly:
\begin{equation}
                         (R_f+d_VR_0)|_{0}=M.
 \label{ch17:flat-current-map}
\end{equation}
The smooth linear averaging operator, rather than a rough row along the root contours,
therefore controls its equal-order Taylor terms. All derivatives of
$R_f+d_VR_0$ have finite $C_{j,L}$ bounds on the common coarse strip.
Their fixed contours have length $O_j(L)$ and remain within bounded
coarse distance; differentiation adds only finite contour-length factors.

For the exact action at the critical configuration the current identity is preserved before
any estimates. The envelope theorem differentiates the stationary value
only through the explicit coarse field. In the penalty term, simultaneous
left multiplication of the averaging operation and the coarse link leaves the
penalty term unchanged. Thus a fine vector field $Y$ induces the coarse
vector field $R_fY$; its root correction $d_VR_0Y$ has zero action
variation by coarse gauge invariance. Together with
\eqref{ch17:block-algebra} this gives
\[
 DS_{\rm stat}[Y'_b]
  =L^5\mathcal I K_0DS[Y_b],\qquad
 \sum_a\nabla_a^{\rm c}J'_{ab}
  =L^5\mathcal I K_0\sum_a\nabla_aJ_{ab}.
\]
Joint stationarity in the fine field and the local gauge transformation is used here, as in the gauge-fixed
quotient calculation of Chapter~\ref{ch07:matching}. Matching replaces
$S_{\rm stat}$ by an action with identical Taylor coefficients through
field order four. Since $Y'(0)=0$, its action variation still agrees
through order four. Therefore the displayed defect pullback is exact
at every field order used in \eqref{ch17:unnormalized-defect}.
A separate additive $C/(\log\widehat\beta)^2$ matching error is neither
needed nor permitted in this argument.

\subsection{Uniform current and difference-primitive bounds}

Recall that $J_{ab}$ is the lattice stress current and $Y_b$ its
associated gauge-covariant vector field. The functions $q_a$ and
$t_{abc}$ are the forward-difference primitives of its trace and
antisymmetric part, respectively, as defined in
Section~\ref{ch17:current-suppliers}. They are distinct from the
small-field threshold $t$ used in the norms with insertions.

For completeness we describe the accompanying bounds needed in the
regulator criterion. Subtract fixed local reference stress and curvature
vector fields, whose leading Taylor coefficients by scaling dimension are those used initially.
The reference functions have finite stencils and fixed coefficient bounds.\mgubi{Issue 8: gauge invariance of the reference stress is not stated; the division by two curls below uses the Ward identity on each leg.}
The lowest jets of the transferred functions agree with them exactly.
For the stress this follows from the exact rescaling at the end of the proof of preservation of the order of vanishing,
applied with $n=2$ to the Hessian $B$ of $K_aJ_{ab}$. Its values on pairs of polynomial fields with $r+r'<2$ vanish: one leg is then constant,
a constant field is the gradient of a linear potential, and the truncation argument given there, using the exponential decay of the kernel, extends the Ward identity to such potentials. If
$r+r'=2$, the rescaling gives $D^2\mathcal I(K_aJ_{ab})(0)[P,P']=L^{-4}B(P,P')$, which the factor $L^4$ in $J'_{ab}$ compensates.
Moreover $K_a$ is a finite combination of fine translates with coefficient sum one, and a translate of a homogeneous polynomial
differs from it by terms of lower degree, on which the Hessian of $J_{ab}$ vanishes; so $B(P,P')$ is the corresponding value for $J_{ab}$.
Consequently the quadratic coefficients of scaling dimension below four remain zero, and the quadratic
coefficient of scaling dimension four is unchanged.

For the vector field one additionally uses the column sum of the linearized averaging operation.
By \eqref{ch17:flat-current-map}, the corrected differential is exactly
$M$ at zero. The alias formula
for $M$ annihilates every nonzero alias at zero coarse momentum, so its
fine column sum in a fixed direction is constant. Its row sum is $L$;
there are $L^4$ positions within a block. The column sum is therefore $L^{-3}$.
This is an identity for the corrected differential, with no residual
root-gradient term. Since the leading fine curvature of $W_a$ is $L^{-2}$ times
the coarse curvature, the factors in the transfer of the current multiply to
$L^5L^{-3}L^{-2}=1$. This preserves the linear curvature jet of $Y$.
Gauge covariance then fixes its accompanying quadratic commutator term:
the covariant polynomial of scaling dimension two in the transferred
vector field is again the curvature.

The stress remainder's quadratic symbol now has order at least three.
Two-leg curl division leaves one extra difference, so its diagonal
transfer is $C_jL^4L^{-5}\le C_j/L$. Higher stress orders have
radial-gauge bounds $C_jL^{4-2r}$.
For $Y$, its zero-order corrected differential $M$ has row sum at most
$C_jL$, while $\|K_0\|_\infty\le C_jL^{-4}$.
To see the resulting local gain, write a translated current coefficient
as $y(f-x;z_1-x,\ldots,z_r-x)$, with output link $f$ and base point $x$.
For each fixed coarse output edge $e$ its absolute sum is bounded by
\[
 \|K_0\|_\infty\sum_f|M_{ef}|
        \sum_x|y(f-x;\,\cdot\,)|
                   \le C_jL^{-3}\|y\|.
\]
Only a bounded number of coarse output edges are within the patches of the averaging
operation and of the base-point average; exponential coefficient tails control the rest.
This proves the same gain for the full local output-link norm.
The linear vector-field remainder has a curvature difference and one
additional difference. If the extra difference acts on its field leg,
it contributes a factor $L^{-1}$ through $W_a$. If it acts on the fine output index,
summation by parts moves it onto $M_{ef}$; the identity for the linear averaging operator gives
$\sum_f|\nabla_fM_{ef}|\le C_j$, again gaining $L^{-1}$ relative to
its row sum $C_jL$. Thus its diagonal factor is
$C_jL^5L^{-3}L^{-3}=C_j/L$.
For field order $r\ge2$ its diagonal factor is $C_jL^{2-2r}$.
Derivatives of the averaging operation or of the correction from its local gauge transformation multiply lower
field orders of $Y$, so their fixed-$L$ constants are absorbed by the
same logarithmic normalization as above. Evaluating in a local radial
gauge is legitimate for this covariant vector field: transform the local gauge
transformation at its output and the differential of the averaging operation together. The additional
derivatives of the local gauge transformation again occur only below the field-order diagonal.

The trace and antisymmetry primitives start at field order two. Their
old terms carry $L^3$ and hence gain $C_jL^{3-2r}\le C_j/L$.
Their new terms are the explicit $L^4K_a^{\rm t}J$ terms proved in
Lemma~\ref{ch17:block-identities}; the equal-order inhomogeneous term from $J$
is bounded independently of $L$. Choose a fixed smaller weight for
these primitive norms so this one-way inhomogeneous term has norm at most $1/8$.
The other diagonal terms are then made small by $L$, and all lower-order
terms by the lower bound on the reference coupling. Reference functions contribute finite
$C_{j,L}$ inhomogeneous terms, divided by positive powers of $\log\widehat\beta$;
a sufficiently large lower bound on the reference coupling can make these terms at most $1/4$.

Thus a weighted sum of the normalized current and primitive norms obeys
$X_m\le\frac12X_{m+1}+\frac14$. The initial Wilson current and zero
primitives have bounded normalized norms. Uniform bounds follow throughout
the history, and in particular at each fixed scale. At that scale
these are the unnormalized finite constants required by the regulator
criterion, since $\widehat\beta_m$ is fixed. The needed free scaling-dimension
jets are unchanged, and the defect tends to zero by the homogeneous
contraction above. All estimates use the same averaging operation, critical
configuration, action matching and periodized kernels as the trajectory with insertions.
